\documentclass[11pt]{article}
\usepackage{amsmath,amssymb,hyperref}
\begin{document}

\section{Outcome and the common-configuration requirement}

This note continues the Fusion project by constructing large exact vacuum-field families, giving two geometric reformulations, and proving new ideal-MHD and particle-confinement estimates. It does not establish a complete hot, stable stellarator with self-consistent collisional kinetics and physical sources and boundaries. The distinction matters because the positive results below currently use different configurations and different evolution equations.

The algebraic family starts from the project's exact Dommaschk vacuum field. The original exact weakly quasisymmetric field is a different seed, with the previously established equilibrium obstruction. The weak-QS particle theorem cannot simply be attached to an arbitrary member of the vacuum family. The stable free-boundary equilibrium constructed below is on a periodic cylinder, whose Euclidean torus extension remains a specific unfinished theorem.

The main completed results are as follows.
\begin{enumerate}
\item The exact Dommaschk potential is a real Laurent polynomial on a complex algebraic threefold. A null-conic construction supplies all homogeneous harmonic polynomials, and hence arbitrarily large finite-dimensional families of exact vacuum perturbations. A KAM argument preserves two nested invariant tori for every sufficiently small coefficient vector, using the nonzero normal twist in the existing project source. The coefficient threshold has not been numerically bounded here.
\item A vacuum field has an exact Abelian Hermitian--Yang--Mills lift to a complex surface with the physical product metric. A separate Euclidean log-determinant quantity is positive and subharmonic, with quantitative derivative bounds. Both identities are static and have explicitly stated boundary conditions.
\item A periodic cylindrical plasma has positive pressure and a strictly positive complete ideal-MHD interface energy for every Fourier mode in an explicit parameter window. The estimate includes the unfavorable curvature term and the zero mode, and is uniform in axial period. A different field branch has an explicit unstable resonant mode.
\item For weak QS, an explicit quadratic-in-velocity correction cancels the first full-orbit momentum defect. This gives a finite-time hot-core theorem with a genuinely absorbing boundary and Maxwellian tails, conditional on an explicit geometric buffer inequality. The orbit identity has an exact relativistic extension. Self-consistent fields, Coulomb collisions and heating are additional terms, not hidden assumptions.
\end{enumerate}

\paragraph{Verification status.}
The new mathematical arguments are written out below and have been checked algebraically and by independent cross-review. Exact rational symbolic calculations verify the converted vacuum potential and its symmetries. A separate rational certificate verifies the numerical Maxwellian lower bounds. These checks are not formal proofs of the whole research note.

For the seed tori, the pinned project source is commit \texttt{6312f1cc9c30b32798ad7b7540766152ef03e76f}. All 227 project modules transitively imported by its nesting theorem were retrieved and matched their Git blob hashes. The source audit found no project axiom declarations, \texttt{sorry}, \texttt{admit}, \texttt{mkSorry}, \texttt{ofReduceBool}, or actual use of \texttt{native\_decide} in that closure. Final theorem assemblies discharge the intermediate residual hypotheses. Witness-evaluation code was distinguished from kernel-checked constructor proof terms. The Lean build and a kernel axiom printout were not rerun here; external Lean and Mathlib dependencies were not independently audited. Consequently the geometric family theorem explicitly uses the reported source theorems as inputs.

\section{The auxiliary-PDE idea and the sign it must preserve}

The relevant part of the OpenAI manuscript is a fiberwise quaternionic Monge--Ampère equation coupled to a positive conserved lifted wave current $A$. Its potential satisfies a prescribed fiber determinant. A corrected fiber energy then obeys an identity of the form
\[
\Delta_{(T,\tau)}F+\int A:\Sigma
=\text{controlled boundary flux},\qquad \Sigma\ge0.
\]
The determinant/cofactor calculation cancels a term without a favorable sign. The resulting inequality is elliptic in physical time and two auxiliary real coordinates; it is not a monotone functional of physical time alone. Sections 2--4 and Proposition 4.2 were examined for this mechanism, without claiming a complete audit of the manuscript or its formalization.\cite{wave}

For plasma, the analogous construction must respect the physical momentum balance and the actual loss boundary. A positive particle moment tensor is available, but it has a Lorentz-force divergence. Adding Maxwell stress restores a total conservation law while introducing an indefinite spatial stress. Choosing an auxiliary determinant does not by itself reconcile these two properties. Likewise, the determinant of a phase-space flow is insufficient: a volume-preserving flow can have determinant one while stretching strongly in one direction and contracting in another.

The calculations below pursue three concrete alternatives: an exact gauge lift of the vacuum equations, a determinant of the physical Hessian, and a corrected particle invariant together with covariant derivatives. Their equations are explicit enough to identify which terms a further auxiliary-PDE construction must control. The cylindrical calculation adds a global positive boundary operator that could serve as the starting point for a geometric stability argument.

\section{Exact algebraic vacuum families and moving confinement domains}\label{exact-algebraic-vacuum-families-and-moving-confinement-domains}

The following construction produces arbitrary-dimensional families of exact vacuum fields with nearby invariant magnetic boundaries. It also supplies new inputs to the positive-pressure weak MHD construction. The algebraic identities were independently checked in exact arithmetic. The seed torus theorems were inspected at source level; their Lean proofs were not independently rebuilt. No hot kinetic stability conclusion is drawn from these geometric statements.

\subsection{The actual complexification and its real structure}\label{the-actual-complexification-and-its-real-structure}

Put

\[
u=x+iy,\quad v=x-iy,\qquad
X=(\mathbb C^*)^2\times\mathbb C_z,
\qquad q=du\,dv+dz^2.
\]

Here the symmetric bilinear metric has entries \(q_{uv}=q_{vu}=1/2\). Its volume form, compatible with the usual orientation on the real locus, is \((i/2)du\wedge dv\wedge dz\), and

\[
\Delta_q=4\partial_u\partial_v+\partial_z^2,
\qquad \nabla_q F=2F_v\partial_u+2F_u\partial_v+F_z\partial_z.
\]

The antiholomorphic involution is

\[
\rho(u,v,z)=(\overline v,\overline u,\overline z).
\]

Its fixed locus is Euclidean three-space with the vertical axis removed. In particular, the real form is not obtained by taking all three complex coordinates to be real.

The locally defined angular potential and its global differential are

\[
\vartheta=\frac{\log u-\log v}{2i},\qquad
d\vartheta=\frac1{2i}\left(\frac{du}{u}-\frac{dv}{v}\right).
\]

The latter is an algebraic regular one-form on \(X\); its meromorphic extension across the deleted coordinate divisors is logarithmic. It is closed and co-closed for \(q\). Its period along a positive real toroidal longitude is \(2\pi\). In algebraic de Rham cohomology its class is \(([d\log u]-[d\log v])/(2i)\).

Thus the precise algebraic class of our fields is

\[
\alpha_G=G\,d\vartheta+dF,\qquad
F\in\mathbb C[u^{\pm1},v^{\pm1},z],\quad \Delta_qF=0.
\]

The physical field is the restriction of \(\alpha_G^{\sharp_q}\) to the fixed locus. It is real when \(G\in\mathbb R\) and \(F=F^*\), where

\[
F^*(u,v,z)=\overline{F(\overline v,\overline u,\overline z)}.
\]

For Laurent coefficients this says \(c_{pqj}=\overline{c_{qpj}}\). Closedness gives curl zero; co-closedness gives divergence zero. Multiplication does not preserve this solution space: \(u\) and \(v\) are harmonic but \(uv\) has Laplacian four.

For the actual Dommaschk candidate, every active radial term has zero logarithmic power. The exact conversion is

\[
R^p\cos(m\phi)=\frac{u^{(p+m)/2}v^{(p-m)/2}+u^{(p-m)/2}v^{(p+m)/2}}2,
\] \[
R^p\sin(m\phi)=\frac{u^{(p+m)/2}v^{(p-m)/2}-u^{(p-m)/2}v^{(p+m)/2}}{2i}.
\]

All exponents are integers. The checked seed has 33 nonzero rational Dommaschk coefficients and 121 nonzero Laurent monomials after collection. The resulting \(F_0\) belongs to \(\mathbb Q(i)[u^{\pm1},v^{\pm1},z]\). Exact checks give

\[
\Delta_qF_0=0,\quad F_0^*=F_0,\quad
F_0(-u,-v,z)=F_0(u,v,z),\quad
F_0(v,u,-z)=-F_0(u,v,z).
\]

The last two identities encode two field periods and stellarator symmetry. The script \texttt{complex\_vacuum\_exact\_check.py} reproduces these algebraic checks; \texttt{dommaschk\_complex\_laurent.json} contains all 121 rational complex coefficients. The script also parses the pinned Lean rational potential data and verifies its exact equality to this Laurent polynomial. These checks are independent of the numerical optimiser metrics and do not certify the invariant tori by themselves.

The Laurent coefficient of \(u^{-6}\) in \(F_0\) is
\[
-\frac{12674582344287}{143360000000000}\,i.
\]
It is nonzero and cannot be cancelled by adding Cartesian polynomials. Consequently the polynomial families below retain nonaxisymmetry about the stated toroidal axis.


\subsection{A complete algebraic harmonic family from the null conic}\label{a-complete-algebraic-harmonic-family-from-the-null-conic}

For \(\lambda\in\mathbb C\), define

\[
Q_\lambda=u+2\lambda z-\lambda^2v.
\]

Its differential is null: \(4(Q_\lambda)_u(Q_\lambda)_v+(Q_\lambda)_z^2=0\). Consequently every holomorphic function \(f(Q_\lambda)\) is harmonic wherever defined. Polynomial powers are entire; rational powers \((Q_\lambda-a)^{-m}\) are admissible on any plasma neighborhood avoiding their pole divisor. Real potentials are obtained by taking \((F+F^*)/2\) and \((F-F^*)/(2i)\).

The map

\[
\lambda\longmapsto[1-\lambda^2:i(1+\lambda^2):2\lambda]
\]

parametrizes the null conic in the projective space of complex covectors. A point of complexified physical space determines the quadratic section \(Q_\lambda\) of \(\mathcal O(2)\to\mathbb{CP}^1\). This is the elementary minitwistor incidence construction \cite{hitchin}. The reality map on the corresponding total space is

\[
(\lambda,\eta)\longmapsto
\left(-\frac1{\overline\lambda},-\frac{\overline\eta}{\overline\lambda^2}\right),
\]

as follows directly from \(\overline{Q_\lambda(\rho p)}=-\overline\lambda^2Q_{-1/\overline\lambda}(p)\). On real space \(Q_\lambda=a\) is a line of direction

\[
n_\lambda=\frac{(-2\Re\lambda,-2\Im\lambda,1-|\lambda|^2)}{1+|\lambda|^2}.
\]

All directions occur. This supplies direction-dependent harmonic building blocks without assuming any estimate from Kakeya theory. Physical Euclidean covectors have no nonzero real null directions; complexification is essential to this construction.

\textbf{Proposition 1 (complete polynomial basis).} Write

\[
Q_\lambda^n=\sum_{k=0}^{2n}H_{n,k}(u,v,z)\lambda^k.
\]

Then the \(2n+1\) functions \(H_{n,k}\) form a complex basis of homogeneous harmonic polynomials of degree \(n\). Real harmonic polynomials of degrees at most \(N\), modulo constants, have dimension \(N(N+2)\).

\textbf{Proof.} Every coefficient is harmonic. Under \((u,v,z)\mapsto(tu,t^{-1}v,z)\), coefficient \(H_{n,k}\) has weight \(n-k\); each coefficient is nonzero, so the coefficients are independent. Conversely, a homogeneous solution is determined by its Cauchy data \(P(u,v,0)\) and \(P_z(u,v,0)\), having respectively \(n+1\) and \(n\) coefficients. The equation recursively determines higher powers of \(z\). Therefore the dimension is at most \(2n+1\), proving completeness. The real fixed subspace has this real dimension; summing over \(n\) gives \((N+1)^2\), before quotienting constants. \(\square\)

There is also a complete basis for the Laurent-polynomial class. Let \((p)_j=p(p-1)\cdots(p-j+1)\). For integers \(p,q\) not both negative, let \(M(p,q)\) be the smallest nonnegative member of \(\{p,q\}\), and put

\[
E_{pq}=\sum_{j=0}^{M(p,q)}\frac{(-4)^j(p)_j(q)_j}{(2j)!}u^{p-j}v^{q-j}z^{2j},
\] \[
O_{pq}=\sum_{j=0}^{M(p,q)}\frac{(-4)^j(p)_j(q)_j}{(2j+1)!}u^{p-j}v^{q-j}z^{2j+1}.
\]

\textbf{Proposition 2.} The \(E_{pq},O_{pq}\) form a basis of \(\ker\Delta_q\) in \(\mathbb C[u^{\pm1},v^{\pm1},z]\).

\textbf{Proof.} Set \(D=\partial_u\partial_v\). Any solution has the unique formal expansion

\[
F=\sum_{j\ge0}\frac{(-4)^jz^{2j}}{(2j)!}D^jA
 +\sum_{j\ge0}\frac{(-4)^jz^{2j+1}}{(2j+1)!}D^jB,
\]

where \(A=F|_{z=0}\), \(B=F_z|_{z=0}\). The sums terminate exactly when \(D\) is locally nilpotent on both data. A Laurent monomial has this property precisely when at least one exponent is nonnegative. Distinct monomials cannot cancel after applying a common power of \(D\). The stated basis follows from the monomial bases for the two data spaces. \(\square\)

\subsection{Density of the polynomial family}\label{density-of-the-polynomial-family}

\textbf{Proposition 3 (harmonic Runge approximation in the needed form \cite{walsh}).} If \(K\subset\mathbb R^3\) is compact with connected complement and \(h\) is harmonic near \(K\), then harmonic polynomials approximate \(h\) uniformly on \(K\). If \(K\) is the closure of a smooth solid torus and \(h\) is harmonic on a fixed larger neighborhood, approximation holds in every prescribed finite \(C^r(K)\) norm and uniformly on a sufficiently small common complex neighborhood of \(K\).

\textbf{Proof.} A measure \(\mu\) annihilating all harmonic polynomials has Newton potential \(U=G*\mu\) identically zero outside a sufficiently large ball: the convergent multipole expansion there consists of harmonic polynomial moments of \(\mu\). Unique continuation and connectedness of \(\mathbb R^3\setminus K\) imply \(U=0\) throughout the complement. Take a cutoff \(\chi\) equal to one near \(K\), supported where \(h\) is harmonic. Distributionally,

\[
\langle\mu,h\rangle=\langle\Delta U,\chi h\rangle
=\langle U,\Delta(\chi h)\rangle=0,
\]

up to the immaterial sign convention for the fundamental solution. The last integrand is supported away from \(K\). Hahn--Banach and Riesz representation give uniform approximation. Apply the result on a slightly thickened solid torus \(K'\) inside the harmonicity neighborhood; its complement is still connected. Interior elliptic estimates give \(C^r\) approximation on \(K\). Interior analytic derivative estimates, with constants depending only on the distance from \(K\) to \(\partial K'\), give uniform convergence of the complex Taylor extensions on a smaller common neighborhood. \(\square\)

Any vacuum field near a solid torus with the same longitudinal circulation as \(B_0\) differs from \(B_0\) by a harmonic gradient: its dual difference is closed and has zero class in the one-dimensional first cohomology. Thus the polynomial construction is dense in the local fixed-circulation vacuum class after a possible shrink of the analytic neighborhood. This is not a claim that arbitrary fixed-domain boundary conditions are preserved.

\subsection{The seed's normal twist is present in the source certificate}\label{the-seeds-normal-twist-is-present-in-the-source-certificate}

The inspected project files at commit \texttt{6312f1cc9c30b32798ad7b7540766152ef03e76f} assert existence of exact analytic invariant parameterizations A and C. The source-to-geometric torus bridge is described below. The toroidal-angle field-line frequency ratios are

\[
\iota_A=-\frac{19+\sqrt5}{89},\qquad
\iota_C=-\frac{85+\sqrt5}{361}.
\]

They are quadratic irrational and therefore satisfy the required Diophantine conditions. Both tori have nonzero \(B_\phi\), and the project nesting statement places C inside A.

In \texttt{DommaschkKamAContract.lean}, \texttt{ExactTorusA.adm} provides \texttt{AdmBoundsG} for the exact torus; the corresponding C theorem is \texttt{ExactTorusC.adm}. The field \texttt{hτ} of this structure states

\[
\tau\leq\left|\operatorname{avg}T_{\rm alg}\right|.
\]

\texttt{DommaschkKamAConsts.lean} and \texttt{DommaschkKamCConsts.lean} take their constants equal to the rational admissibility bounds. The values are

\[
\tau_A=\frac{6997}{500},\qquad \tau_C=\frac{141623}{10000}.
\]

These are source-specific torsion bounds, not derivatives of rotational transform in an arbitrary engineering flux coordinate.

To identify their meaning, use toroidal angle \(\phi\) as time and write \(X=(RB_R/B_\phi,RB_z/B_\phi)\) in the poloidal plane. The preserved time-dependent density is \(w=B_\phi\). Let \(E=K_\theta\). The source normal frame \(N\) is normalized by \(w\det(E,N)=1\), and \texttt{KamFlowFrameG.torsion\_coord} and the coordinate reconstruction identities give

\[
(\iota\partial_\theta+\partial_\phi)N-DX\,N=T E+S N.
\]

At exact invariance the error terms vanish, \(S=0\), and \(T=T_{\rm alg}\). A variation \(aE+bN\) then satisfies \(b'=0\), \(a'=-Tb\). Solving the Diophantine cohomological equation

\[
(\iota\partial_\theta+\partial_\phi)h=\langle T\rangle-T
\]

and replacing \(N\) by \(N+hE\) makes the torsion constant. The normal linearization of the full-turn return map is therefore

\[
(a,b)\longmapsto(a-2\pi\langle T\rangle b,b).
\]

In the project's phase convention \(\psi=2\phi-\pi/2\), the operator is instead written \(\iota\partial_\theta+2\partial_\psi\); the same conclusion follows. This proves the needed averaged twist is nonzero, using the reported exact torus and admissibility theorems. The certificate itself has not been rebuilt here.

The source delivers invariant parameterizations and poloidal winding data. To obtain embedded tori, first note that positivity of the cylindrical radius and nonsingularity of the certified normal frame give an immersion. Equality of two image points fixes the toroidal angles modulo a full turn. Uniqueness of the field flow and density of the irrational torus orbit then make their poloidal parameter difference a global translation period. An irrational such period forces a constant poloidal loop; a rational period of order \(q>1\) makes its winding divisible by \(q\). Both contradict winding one. The compact injective immersion is therefore an embedding. The certified polar radial separation gives nesting of the resulting bounded solid tori. This geometric argument is supplied here, rather than represented as another rebuilt Lean certificate.

\subsection{A family theorem with moving physical domains}\label{a-family-theorem-with-moving-physical-domains}

\textbf{Theorem 4.} Assume the exact seed torus and admissibility theorems in the inspected project source. For every \(N\geq1\), choose a real basis \(P_1,\ldots,P_{N(N+2)}\) of harmonic polynomials of degree at most \(N\), modulo constants. There is \(\epsilon_N>0\) such that for every real coefficient vector \(|a|<\epsilon_N\),

\[
B_a=B_0+\nabla\left(\sum_j a_jP_j\right)
\]

is an exact analytic vacuum field and has two nearby nested analytic invariant tori \(\Sigma_{C,a}\subset\Omega_{A,a}\), \(\partial\Omega_{A,a}=\Sigma_{A,a}\), with the original rotation ratios \(\iota_C,\iota_A\). It has longitudinal circulation \(2\pi\), zero meridional circulation, and strictly positive toroidal magnetic flux. Field lines beginning in the closed outer solid torus stay there for all positive and negative field-line time. Distinct coefficient vectors produce distinct outer domains. The complex field remains an algebraic logarithmic form plus an exact rational harmonic form.

\textbf{Proof.} Vacuum identities and periods follow immediately. Work on a fixed complex neighborhood of the two seed tori and a real solid-torus neighborhood of their interiors. Coefficient smallness makes the field perturbation small in a fixed analytic norm. Nonvanishing of \(B_\phi\) near each seed torus gives an analytic annular return map, close to the original map.

The return map is \textbf{exact} symplectic, not merely area preserving. Indeed \(\beta_a=\iota_{B_a}\mathrm{vol}\) is closed and extends across the solid torus, whose second cohomology vanishes, so \(\beta_a=dA_a\). If \(P_a\) is a return map, \(T_a(x)\) its return time, and \(\lambda_a=A_a|_{\rm section}\), the homotopy formula gives

\[
P_a^*\lambda_a-\lambda_a=dS_a,\qquad
S_a(x)=\int_0^{T_a(x)}A_a(B_a)(\Phi_t(x))\,dt.
\]

The term involving \(dA_a\) vanishes because \(\iota_{B_a}\beta_a=0\). One may change both signs to use the source area form \(B_\phi\,dR\wedge dz\). A smooth global primitive suffices for this cohomological argument; the analytic annular area form has a local analytic primitive, and its exactness property is the same because the return map acts trivially on the first cohomology of the annulus.

The preceding section verifies nonzero averaged twist. The analytic a posteriori KAM theorem for exact symplectic maps, with prescribed Diophantine rotation \cite{dlL}, now gives a nearby invariant circle for each of the two return maps. The suspension yields analytic embedded tori. Taking perturbations sufficiently small preserves their embedding, separation, nesting, and toroidal winding. This works for every coefficient vector in an open small ball: the rotation numbers are fixed at the two Diophantine seed values while normal positions adjust. No Cantor restriction on the external coefficient vector is needed.

Tangency to the compact outer boundary makes the closed solid torus invariant under the complete restricted field flow. To prove positive toroidal flux, cut the torus along a meridional disk. Integration by parts with a single-valued potential on the cut domain gives

\[
\int_{\Omega_{A,a}}|B_a|^2\,dx=\Gamma_a\Phi_a=2\pi\Phi_a>0.
\]

Finally, if two coefficient vectors gave the same domain, their difference would be the gradient of a harmonic function with zero normal derivative there. Green's identity forces that gradient to vanish. Polynomial independence modulo constants forces the coefficient vectors to coincide. \(\square\)

The same construction works with arbitrary finite collections of the real Laurent or pole-free rational harmonic functions above. An infinite-dimensional analytic version follows by taking an absolutely convergent coefficient series in a weighted analytic norm on a fixed complex ball containing the plasma neighborhood and requiring its norm to be below the common KAM threshold. Infinite sums need not remain rational functions. No numerical lower bound for that threshold has been extracted here.

If an exactly fixed magnetic flux \(\Phi_*\) is required, replace \(B_a\) by \((\Phi_*/\Phi_a)B_a\). This preserves vacuum, all field lines and invariant tori. It varies the longitudinal circulation. Fixed flux and fixed circulation are different requirements; the theorem does not silently identify them. The normalized family still has distinct domains by one-dimensional Hodge uniqueness.

A symmetry-preserving infinite family is available by projecting harmonic polynomials under \(R(x,y,z)=(-x,-y,z)\) and \(S(x,y,z)=(x,-y,-z)\):

\[
P\mapsto\tfrac14(P+P\circ R-P\circ S-P\circ RS).
\]

For an explicit independent family, take

\[
h_m=\operatorname{Im}(x+iy)^{2m}\quad(m\geq1),\qquad
k_m=z\operatorname{Re}(x+iy)^{2m}\quad(m\geq0).
\]

Both families are harmonic and have the required parity. Thus two field periods and stellarator symmetry can be retained without restricting to a finite set of perturbations. The physical vector-field laws are \(B(Rx)=RB(x)\) and \(B(Sx)=-SB(x)\): the latter is a reversing field-line symmetry. In cylindrical components, \(B_R(R,-\phi,-z)=-B_R(R,\phi,z)\), whereas \(B_\phi\) and \(B_z\) are even under the simultaneous reversal.

\subsection{Connection to the positive-pressure construction and its limit}\label{connection-to-the-positive-pressure-construction-and-its-limit}

The weak MHD framework is related to \cite{elps}. The toroidal-time parameterization of each perturbed torus satisfies \(B_a\circ K=f DK\omega\), with \(f=B_\phi/R>0\), \(\omega=(\iota,1)\). Nonvanishing gives a constant sign for \(f\); averaging \(K^*\alpha(\omega)=f|DK\omega|^2\) gives one because the angular circulation is \(2\pi\), proving that the sign is positive. It can be rectified to physical time: put \(m=\langle f^{-1}\rangle\), solve

\[
\omega\cdot\partial h=f^{-1}/m-1,
\qquad F(\theta)=\theta+\omega h(\theta).
\]

The Jacobian determinant of \(F\) is \(f^{-1}/m>0\), so its degree-one map is a diffeomorphism. With \(K'=K\circ F^{-1}\), one gets \(DK'(\omega/m)=B_a\circ K'\). Consequently every member satisfies the analytic Diophantine boundary-flow hypotheses of the previously derived positive-pressure weak MHD lift. That lift uses an invariant graph of the boundary kinetic Hamiltonian, an analytic vacuum Cauchy extension, an edge current sheet and a supported exterior current sheet.

This produces a much larger family of positive-pressure weak MHD equilibria. It does not prove free-boundary nonlinear stability, a smooth finite-beta pressure profile, coil manufacturability, kinetic confinement, collisional heating balance, or a high-norm continuation bound. In particular, invariant magnetic field lines do not trap every finite-Larmor-radius particle trajectory.

\paragraph{Source pin and trust boundary}\label{source-pin-and-trust-boundary}

The source snapshot is \href{https://github.com/siddhartha-gadgil/PlasmaLean/commit/6312f1cc9c30b32798ad7b7540766152ef03e76f}{commit 6312f1cc9c30b32798ad7b7540766152ef03e76f}, dated 17 November 2026. In the \texttt{Design} directory, \texttt{DommaschkKamAFinal.exists\_exactTorusA} closes the residual hypotheses using \texttt{DommaschkKamARes}; \texttt{DommaschkKamCFinal} closes the corresponding C hypotheses. \texttt{DommaschkKamNested.dommaschk\_nested} then constructs both parameterizations and the nesting data. \texttt{ExactTorusA.adm} and \texttt{ExactTorusC.adm} provide the admissibility bounds used above, including nonzero averaged torsion. These are source theorem assertions with supplied proof bodies, not additional uninstantiated hypotheses in the final statements. A successful rebuild and an explicit Lean axiom report remain separate verification steps.

A further dependency audit fetched all 227 project modules in the transitive import closure of \texttt{DommaschkKamNested} at the stated commit and verified their Git blob hashes. No project axiom declaration, \texttt{sorry}, \texttt{admit}, \texttt{mkSorry}, \texttt{ofReduceBool}, or actual use of \texttt{native\_decide} was found. Six \texttt{unsafe} expressions occur in meta-level numerical witness evaluators in \texttt{KronProof} and \texttt{KronProof2D}; the proof builders construct ordinary constructor and reflexivity terms for subsequent kernel checking. The auxiliary \texttt{KernelFact} tactic likewise builds explicit proof terms. This audit supports the stated source trust boundary but is not a substitute for compiling the proof closure or auditing the external Lean and Mathlib dependencies.


The exact original rational candidate is available at \url{https://raw.githubusercontent.com/siddhartha-gadgil/PlasmaLean/6312f1cc9c30b32798ad7b7540766152ef03e76f/scripts/dommaschk_qa/results/20261031T092522Z/candidate.json}. Its optimiser metadata is explicitly labelled untrusted; the algebraic field check and the later proof chain are distinct from that numerical fit.

\section{A complex-surface lift that retains the Euclidean metric}

The algebraic complexification above has complex dimension three. There is a different, exact construction on a complex surface: add one auxiliary circle to physical three-space. This is a static gauge-theoretic lift, not an identification of that circle with plasma evolution time.

Let $\Omega$ be a smooth solid torus, let $B$ be a smooth real field on a neighborhood of its closure, and set
\[
\alpha=B^\flat,\qquad \beta=*_{3}\alpha,\qquad
M=\Omega\times(\mathbb R/\ell\mathbb Z).
\]
Give $M$ the product Euclidean metric and orientation $dx\wedge dy\wedge dz\wedge dt$. Define the real two-form
\[
\mathcal F=\beta-\alpha\wedge dt.
\]

\paragraph{Proposition (exact Abelian Hermitian--Yang--Mills lift).}
The field $B$ is vacuum if and only if $d\mathcal F=0$. The form $\mathcal F$ is anti-self-dual for every $B$. Consequently a vacuum field gives a closed primitive form of type $(1,1)$ for the complex coordinates
\[
u=x+iy,\qquad w=\exp(2\pi(z+it)/\ell).
\]
If $\Omega$ avoids the vertical axis, $M$ is a complex domain in $(\mathbb C^*)^2$. Its metric and Kähler form are
\[
g_4=|du|^2+\frac{\ell^2}{4\pi^2|w|^2}|dw|^2,
\qquad
\omega=\frac i2du\wedge d\bar u+
\frac{i\ell^2}{8\pi^2|w|^2}dw\wedge d\bar w.
\]
In particular, the physical metric is not replaced by the standard flat Hermitian metric in the algebraic coordinates $(u,w)$.

\paragraph{Proof.}
The components of $d\mathcal F=d\beta-d\alpha\wedge dt$ with and without $dt$ vanish separately. Their vanishing says respectively $\operatorname{div}B=0$ and $\operatorname{curl}B=0$. The product Hodge star satisfies
\[
*_{4}\beta=\alpha\wedge dt,\qquad
*_{4}(\alpha\wedge dt)=\beta.
\]
Thus $*_{4}\mathcal F=-\mathcal F$. On this Kähler surface, real anti-self-dual two-forms are precisely the real primitive $(1,1)$ forms. This last assertion can also be verified in the orthonormal coordinates $(x,y,z,t)$; no compactness or algebraic compactification is used. The metric formula follows by differentiating $w$. $\square$

\paragraph{Integral periods and the holomorphic line bundle.}
Let $\gamma$ generate the longitudinal first homology of $\Omega$, and write $\Gamma=\int_\gamma\alpha$. Since $H^2(\Omega;\mathbb R)=0$, the only possible nonzero period of $\mathcal F$ on $M$ is
\[
\frac{1}{2\pi}\int_{\gamma\times S^1_\ell}\mathcal F
=-\frac{\Gamma\ell}{2\pi}.
\]
For the normalized seed, $\Gamma=2\pi$ and $\ell=1$ make this period integral. More generally one may choose the auxiliary period or an overall mathematical curvature normalization. This is not a physical flux-quantization assertion.

Choose a global one-form $a$ on $\Omega$ with $da=\beta$. Locally write $\alpha=d\Phi_i$ and set $A_i=a-\Phi_i\,dt$. Then $dA_i=\mathcal F$. A potential branch change by $2\pi n$ changes $A_i$ by $-2\pi n\,dt$, which has the unitary transition $e^{2\pi int}$ when $\ell=1$. With convention $\nabla=d-iA$, the curvature is $-i\mathcal F$ and $c_1=[\mathcal F/(2\pi)]$. The type $(1,1)$ condition makes the $(0,1)$ part of this connection integrable, hence defines a holomorphic line bundle; the primitive condition is the Abelian Hermitian--Yang--Mills equation with zero Einstein constant. The domain need not be assumed Stein for this explicit construction.

The logarithmic differential on the threefold is algebraic, but this line bundle on the surface need not be algebraic. The algebraic Picard group of $(\mathbb C^*)^2$ is zero: its coordinate ring is a Laurent polynomial unique factorization domain and is regular. A bundle with the nonzero Chern period above therefore cannot be the restriction of an algebraic line bundle on that ambient algebraic torus. This is a genuine distinction between the two constructions, not an inconsistency.

\paragraph{A positive auxiliary energy and its exact scope.}
For a vacuum field,
\[
\int_M|\mathcal F|^2=2\ell\int_\Omega|B|^2.
\]
If $a_1$ is a real one-form whose tangential trace on $\partial M$ vanishes, then $d*\mathcal F=0$ and integration by parts give
\[
\int_M|\mathcal F+da_1|^2-\int_M|\mathcal F|^2
=\int_M|da_1|^2\ge0.
\]
Thus the lift minimizes the Abelian Yang--Mills energy against these fixed-boundary connection variations. A moving plasma boundary, plasma pressure, distribution functions, collisions, and source work are absent from this identity. It is not a plasma Lyapunov theorem.

\section{A determinant quantity for vacuum geometry}

Here is a separate elementary determinant construction that retains the physical Euclidean Hessian. It gives an exact positive identity and useful geometric control, although not a dynamical plasma estimate.

\paragraph{Proposition (a subharmonic log-determinant).}
Let $B$ be a smooth vacuum field on a neighborhood of the closure of a bounded smooth domain. Set $H=\nabla B$, choose $\lambda>0$, and suppose $\|H\|_{\mathrm{op}}\le\delta\lambda$ on the boundary, with $0<\delta<1$. Then throughout the domain
\[
A=I+\lambda^{-1}H>0,\qquad
Q=-\log\det A\ge0,
\]
and
\[
\Delta Q=\sum_k\operatorname{tr}
\big(A^{-1}(\partial_k A)A^{-1}(\partial_k A)\big)\ge0.
\]
Moreover,
\[
\frac{|H|_F^2}{2\lambda^2(1+\delta)^2}
\le Q\le
\frac{|H|_F^2}{2\lambda^2(1-\delta)^2},
\qquad
\Delta Q\ge\frac{|\nabla H|_F^2}{\lambda^2(1+\delta)^2}.
\]

\paragraph{Proof.}
Vacuum implies $H^T=H$, $\operatorname{tr}H=0$, and $\Delta H=0$. For every fixed real vector $\xi$, the scalar $\xi^TH\xi$ is harmonic. The maximum principle therefore extends the boundary spectral bounds $-\delta\lambda I\le H\le\delta\lambda I$ throughout the domain. Differentiating $-\log\det A$, and using $\Delta A=0$, gives the displayed Laplacian identity. Each summand is the squared Frobenius norm of $A^{-1/2}(\partial_kA)A^{-1/2}$. For the comparison bounds, let $h_j$ be the eigenvalues of $H/\lambda$. Since $\sum_jh_j=0$,
\[
Q=\sum_j\big(h_j-\log(1+h_j)\big).
\]
The scalar function in parentheses has value and first derivative zero at the origin, and second derivative $(1+h)^{-2}$. Taylor's integral formula on $[-\delta,\delta]$ gives the two-sided bounds. The lower bound for $\Delta Q$ follows by the same spectral bounds on $A^{-1/2}$. $\square$

Locally, $B=\nabla\Phi$ and $A$ is the Hessian of $|x|^2/2+\Phi/\lambda$. Different branches of the angular potential differ by constants, so $A$ and $Q$ remain globally well-defined. This is a real Monge--Ampère-type quantity with a fully checked sign. The maximum principle controls $Q$ from its boundary data. It does not require, or assert, a globally single-valued convex potential on the solid torus.

The limitations are explicit. A general MHD or kinetic perturbation has current; then $\nabla B$ need not be symmetric and $\Delta\nabla B$ need not vanish. The displayed identity does not close a plasma evolution estimate. A successful auxiliary-PDE argument would have to include those terms and the physical boundary terms with controlled signs.

\paragraph{A tempting simplification is impossible.}
A nonzero smooth vacuum field tangent to the boundary of a bounded smooth domain cannot have constant strength on the entire boundary. To prove this, put $q=|B|^2$. Then $\Delta q=2|\nabla B|^2\ge0$. If $q=c>0$ on the boundary and is nonconstant inside, the maximum principle and Hopf lemma give $\partial_nq>0$ at every boundary point. Tangency and curl-freeness instead give
\[
\frac12\partial_nq=-\mathrm{II}(B,B),
\]
where $\mathrm{II}(v,w)=\langle D_vn,w\rangle$. At a point maximizing distance from a fixed interior point, comparison with the enclosing sphere gives $\mathrm{II}\ge R^{-1}g$ on the tangent plane. There $B\ne0$, so the last normal derivative is negative, a contradiction. If $q$ is constant throughout, the Bochner identity forces $\nabla B=0$; a nonzero constant vector cannot be tangent to the boundary of a bounded domain, as integrating $\operatorname{div}((B\cdot x)B)=|B|^2$ shows. The case $c=0$ forces $B=0$ by the maximum principle. This excludes that particular attempt to remove field-strength variation; it is not an instability theorem.

\section{A full positive energy calculation in a cylindrical benchmark}

The following calculation signs the entire ideal-MHD quadratic energy,
including the interface curvature term.  It is a benchmark on a periodic
Euclidean cylinder, not a claim that a periodic cylinder is an embedded
stellarator in three-dimensional Euclidean space.  Its useful feature is
that the estimates are uniform in the axial period.

Let the plasma occupy the solid periodic cylinder (including its axis,
where all fields and displacements are required to be smooth)
\[
 r<a,\qquad \theta\in\mathbb R/(2\pi\mathbb Z),
 \qquad z\in\mathbb R/(L\mathbb Z),
\]
and let the vacuum occupy $a<r<b$.  The outer surface $r=b$ is fixed and
perfectly conducting.  Take the equilibrium fields
\[
 Y=B_i e_z,\qquad
 X=B_e e_z+B_\theta\frac{a}{r}e_\theta,
 \qquad B_i>0.
\]
Both fields are curl-free and divergence-free in their respective open
regions.  The plasma pressure and density are positive constants.  The
interface has the surface current required by the jump of the tangential
magnetic field.  Force balance is
\[
 p=\frac{B_e^2+B_\theta^2-B_i^2}{2\mu_0}.
\]

We use the lossless ideal-MHD energy principle.  In particular, the outer
wall is passive, there is no work from a perturbed external power supply,
and the magnetic fluxes through the independent material cuts are fixed.
One convenient complete choice fixes the plasma axial flux, the vacuum
axial flux, and the vacuum azimuthal flux.  The last two conditions specify
the harmonic parts of the vacuum perturbation.  The condition that the
outer wall is fixed is not, by itself, a specification of these harmonic
parts.  The calculations below also remain lower bounds if additional
admissible harmonic vacuum perturbations are retained: their magnetic
energy is nonnegative and orthogonal to the gradient part.

Write $\eta=\xi\cdot e_r$ for the interface displacement and
\[
 Q=\nabla\times(\xi\times Y)
\]
for the plasma magnetic perturbation.  If $b_v$ is the vacuum magnetic
perturbation, the complete quadratic energy is
\[
\begin{split}
2\delta W={}&
 \int_{r<a}\left(\frac{|Q|^2}{\mu_0}
       +\gamma p\,|\operatorname{div}\xi|^2\right)\,dx
 +\frac{1}{\mu_0}\int_{a<r<b}|b_v|^2\,dx
 -\frac{B_\theta^2}{\mu_0 a}\int_{r=a}|\eta|^2\,dS.
\end{split}
\]
Here $\gamma>1$ is the adiabatic exponent.  In particular, the negative
surface term has not been omitted.  Its coefficient follows from
\[
 C=\frac{\mathrm{II}(Y,Y)-\mathrm{II}(X,X)}{\mu_0}
   =-\frac{B_\theta^2}{\mu_0a},
\]
where $\mathrm{II}(v,w)=v\cdot\nabla_w n$ and $n=e_r$.

\subsection{Exact mode minimization}

Expand
\[
 \eta(\theta,z)=\sum_{m,n\in\mathbb Z}
       \eta_{mn}\exp(i m\theta+i k_nz),
 \qquad k_n=\frac{2\pi n}{L}.
\]
The normalization is
\[
 \frac{1}{2\pi aL}\int_{r=a}|\eta|^2\,dS
   =\sum_{m,n}|\eta_{mn}|^2.
\]
Put $\kappa=a|k|$.  Derivatives of $I_m$ and $K_m$ below mean derivatives
with respect to their arguments.

First suppose $k\ne0$, and replace $m$ by $|m|$ when defining the radial
functions.  If $w$ is the transverse part of $\xi$, then
\[
 Q_\perp=i kB_iw,\qquad Q_z=-B_i\operatorname{div}_\perp w.
\]
Consequently
\[
 \int |Q|^2\,dx
  =B_i^2\int\left(k^2|w|^2+
                         |\operatorname{div}_\perp w|^2\right)\,dx.
\]
The minimum for prescribed $w\cdot e_r=\eta$ is attained by
$w=\nabla_\perp\chi$, where
\[
 (\Delta_\perp-k^2)\chi=0,
 \qquad \partial_r\chi(a)=\eta.
\]
Indeed, writing $w=\nabla_\perp\chi+v$ with $v\cdot e_r=0$, the cross term
vanishes by integration by parts.  The minimizing longitudinal
displacement is $\xi_z=i k\chi$, which gives
$\operatorname{div}\xi=0$.  Thus dropping the positive compressibility
term loses nothing for these nonzero-$k$ modes.  The plasma contribution
per unit interface area and per $|\eta|^2$ is
\[
 \frac{B_i^2}{\mu_0a}
       \frac{\kappa I_m(\kappa)}{I_m'(\kappa)}.
\]

At the inner boundary the vacuum condition is
\[
 b_v\cdot e_r=\operatorname{div}_{\Sigma}(\eta X)
       =i\left(kB_e+\frac{mB_\theta}{a}\right)\eta,
\]
and $b_v\cdot e_r=0$ at $r=b$.  For a nonconstant Fourier mode, the vacuum
perturbation is the gradient of a single-valued harmonic potential.  Let
$N_{\rm out}(m,k)>0$ be its Neumann impedance, so that its energy per unit
inner-boundary area is
\[
 N_{\rm out}(m,k)
       \left|kB_e+\frac{mB_\theta}{a}\right|^2|\eta|^2.
\]
Set $n_{\rm out}=N_{\rm out}/a$.  The exact minimized coefficient is
\[
 \frac{2\delta W_{mn}}
      {(2\pi aL)|\eta_{mn}|^2}
 =\frac{B_i^2}{\mu_0a}\,F_{mn},
\]
where, for $k\ne0$,
\[
 F_{mn}=
 \frac{\kappa I_m(\kappa)}{I_m'(\kappa)}
 +\left(\frac{k aB_e+mB_\theta}{B_i}\right)^2 n_{\rm out}
 -\frac{B_\theta^2}{B_i^2}.
\]
For arbitrary plasma displacements the right side is a lower bound;
the omitted difference is nonnegative.

For $k=0$ and $m\ne0$, a transverse divergence-free displacement gives
zero plasma magnetic and compression energy.  Direct solution of the
vacuum problem gives
\[
 N_{\rm out}(m,0)=\frac{a}{|m|}
        \coth\left(|m|\log\frac ba\right),
\]
and hence
\[
 F_{m0}=\frac{B_\theta^2}{B_i^2}
     \left(|m|\coth\left(|m|\log\frac ba\right)-1\right)>0.
\]
In an unbounded exterior this becomes
$(|m|-1)B_\theta^2/B_i^2$: the negative curvature term cancels exactly
for a rigid transverse translation, $|m|=1$.  Thus a negative curvature
coefficient alone is not an instability proof.  The fixed close wall
removes this particular zero mode.

\subsection{An exact unstable branch, even with a close wall}

The choice of pressure-lifting direction matters.  Suppose $B_e\ge B_i$
and $B_\theta\ne0$.  For $m=1$, choose
\[
 k=-\frac{B_\theta}{aB_e}.
\]
The vacuum normal perturbation then vanishes identically, independently
of the wall distance.  With $\kappa=|B_\theta|/B_e$, the identity
$\kappa I_1'=I_1+\kappa I_2$ gives
\[
 \frac{\kappa I_1}{I_1'}<\kappa^2.
\]
Consequently the exact energy coefficient is strictly less than
\[
 \frac{B_\theta^2}{B_e^2}
       -\frac{B_\theta^2}{B_i^2}\le0.
\]
This is a genuine negative-energy mode whenever the indicated $k$ is
allowed by the axial period.  By continuity, every sufficiently long
period admits a nearby unstable mode.  Thus positive pressure,
noncollinearity, and a close wall do not stabilize arbitrary choices
of the two tangential fields.  The stable family below uses a slightly
smaller exterior axial field, while keeping the total exterior magnetic
pressure larger than the interior value.

\subsection{A low-pressure, all-mode stability window}

Here is one explicit sufficient window.  Choose
\[
 0<s\le\frac14,\qquad
 B_\theta^2=sB_i^2,\qquad
 B_e^2=\left(1-\frac{3s}{4}\right)B_i^2.
\]
Both axial fields are nonzero, and
\[
 p=\frac{sB_i^2}{8\mu_0}>0,
 \qquad \frac{2\mu_0p}{B_i^2}=\frac{s}{4}.
\]
Let
\[
 h=\frac ba,\qquad \ell=\log h,
 \qquad D=\frac{h^2-1}{2},
 \qquad \beta_e=|B_e|/B_i=\sqrt{1-\frac{3s}{4}}.
\]
Assume
\[
 \ell+\frac{sD}{\beta_e^2}\le\frac{s}{8}.
\]
A simple sufficient choice is
\[
 1<h\le1+\frac{s}{16}.
\]
This is a close ideal conducting wall, not a distant coil condition.

For $m\ge1$, the positive series for $I_m$ gives
\[
 \frac{I_{m+1}(\kappa)}{I_m(\kappa)}
       \le\frac{\kappa}{2(m+1)}.
\]
Using $\kappa I_m'=mI_m+\kappa I_{m+1}$, we obtain
\[
 \frac{\kappa I_m}{I_m'}
 \ge \frac{\kappa^2}
             {m+\kappa^2/(2(m+1))}.
\]
There is also an elementary lower bound for the vacuum impedance:
\[
 n_{\rm out}\ge\frac{1}{m^2\ell+\kappa^2D}.
\]
To see this, use a radially constant potential as a trial function in
the dual variational characterization of the Neumann energy.  Its
energy per unit inner-boundary area, divided by the squared boundary
value, is $(m^2\ell+\kappa^2D)/a$.  Optimizing its amplitude proves the
displayed lower bound.

The least favorable helicity has vacuum factor
$(m\sqrt{s}-\kappa\beta_e)^2$.  Put
\[
 t=\frac{\kappa\beta_e}{m\sqrt{s}}\ge0.
\]
The preceding estimates give
\[
 \frac{F_{mn}}{s}\ge
 \frac{t^2}{\beta_e^2+st^2/2}
 +\frac{(1-t)^2}{\ell+st^2D/\beta_e^2}-1.
\]
For $0\le t\le1$, the elementary Cauchy inequality
\[
 \frac{t^2}{A}+\frac{(1-t)^2}{B}\ge\frac{1}{A+B}
 \qquad (A,B>0)
\]
and the wall assumption imply
\[
 \frac{F_{mn}}{s}\ge
 \frac{1}{1-s/8}-1\ge\frac{s}{8}.
\]
For $t\ge1$, the first term alone yields
\[
 \frac{F_{mn}}{s}\ge
 \frac{1}{\beta_e^2+s/2}-1
 =\frac{1}{1-s/4}-1\ge\frac{s}{4}.
\]
Therefore
\[
 F_{mn}\ge\frac{s^2}{8}
 \qquad\text{for every }m\ne0\text{ and every }k.
\]
The $k=0$ case can either be read directly from the exact formula or
obtained as the limit of this estimate.

For $m=0$, $k\ne0$, the series inequality $I_1\le\kappa I_0/2$ gives
\[
 F_{0n}\ge\frac{\kappa I_0}{I_1}-s\ge2-s>0.
\]
The vacuum term is nonnegative and has not been needed for this bound.

For completeness, the stated simple wall choice implies the wall
assumption.  Indeed, $\ell\le s/16$, $\beta_e^2\ge13/16$, and
\[
 D\le\frac{s}{16}+\frac{s^2}{512}
     \le\frac{129s}{2048}.
\]
It follows that
\[
 \ell+\frac{sD}{\beta_e^2}
 \le\left(\frac1{16}+\frac{129}{6656}\right)s
 <\frac{s}{8}.
\]

\subsection{The spatially constant displacement and the flux constraints}

For $(m,k)=(0,0)$, it is important not to discard the plasma flux
constraint.  A uniform displacement $\eta$ requires
\[
 \delta B_i=-\frac{2B_i\eta}{a},
 \qquad \operatorname{div}\xi=\frac{2\eta}{a}
\]
for the energy-minimizing radial displacement.  Thus the plasma
magnetic and gas contributions per unit interface area are
\[
 \frac{2B_i^2}{\mu_0a}|\eta|^2
       +\frac{2\gamma p}{a}|\eta|^2.
\]
Even if the nonnegative vacuum contribution is discarded, the full
coefficient satisfies
\[
 F_{00}\ge2-s+\frac{\gamma s}{4}>0.
\]

The harmonic vacuum terms can also be written explicitly.  Fixing the
vacuum axial flux and azimuthal flux gives
\[
 \delta B_e=\frac{B_e\eta}{aD},
 \qquad
 \delta(aB_\theta)=\frac{B_\theta\eta}{\ell}.
\]
Their energies add respectively $\beta_e^2/D$ and $s/\ell$ to
$F_{00}$.  Thus in this flux ensemble
\[
 F_{00}=2-s+\frac{\gamma s}{4}
                +\frac{\beta_e^2}{D}+\frac{s}{\ell}.
\]
These contributions are orthogonal to the nonconstant potential
modes.  A calculation that specified only the outer wall and silently
set all magnetic periods to zero would miss these constraints.

\subsection{An explicit uniform interface gap}

The preceding positive constant bound can be strengthened to a
half-derivative interface estimate.  For $m\ge1$, the integral
representations of the Bessel functions imply
\[
 \frac{I_{m+1}(\kappa)}{I_m(\kappa)}<1,
 \qquad \frac{K_{m-1}(\kappa)}{K_m(\kappa)}\le1.
\]
The first ratio is the expectation of a number in $(-1,1)$ in the
positive integral representation of $I_m$; the second follows from
$\cosh((m-1)t)\le\cosh(mt)$ in the integral representation of $K_m$.
They give
\[
 \frac{\kappa I_m}{I_m'}\ge\frac{\kappa^2}{m+\kappa},
 \qquad n_{\rm out}\ge\frac{1}{m+\kappa}.
\]
For the second inequality, the wall impedance is at least the
unbounded-exterior impedance.  This follows by restricting trial
functions in the dual Neumann variational principle to the smaller
annulus.  The unbounded-exterior value is
$-K_m/(\kappa K_m')$, and the displayed ratio bound applies.

The quadratic form
\[
 \kappa^2+(m\sqrt{s}-\kappa\beta_e)^2
\]
has determinant $s$ and trace
$s+1+\beta_e^2=2+s/4\le33/16$.  Its least eigenvalue is therefore at
least $16s/33$.  Consequently, with $x=m+\kappa$,
\[
 F_{mn}\ge\frac{8s}{33}x-s.
\]
If $x\ge9$, this is at least $4sx/33$.  If $x\le9$, the previously
proved estimate $F_{mn}\ge s^2/8$ applies.  Combining the two gives
the convenient conservative bound
\[
 F_{mn}\ge\frac{s^2}{80}(1+m+\kappa)
 \qquad (m\ge1).
\]
For $m=0$, use $I_1<I_0$ in addition to the earlier series bound to
obtain $F_{0n}\ge\max(2-s,\kappa-s)$; this is more than sufficient for
the same estimate.  The zero mode was checked separately above.

We have therefore proved, for every smooth admissible displacement,
\[
 2\delta W\ge
 \frac{2\pi aL\,B_i^2}{\mu_0a}\,
 \frac{s^2}{80}
 \sum_{m,n}\left(1+|m|+a|k_n|\right)|\eta_{mn}|^2.
\]
The constant is independent of $L$.  This is a genuine positive
interface $H^{1/2}$ energy gap for the complete quadratic functional.
The nonnegative bulk excess over the mode minima is still present.
There remain bulk relabelling directions with $\eta=0$, $Q=0$ and
$\operatorname{div}\xi=0$.  The theorem does not give a coercive norm
of every material displacement, a nonlinear continuation theorem, or
a kinetic confinement theorem.

\section{A geometric sufficient condition on an actual toroidal interface}

Use the same passive, lossless exterior and fixed-flux ensemble as above.
Let $Y$ be a tangent vacuum field on an analytic toroidal interface,
straightened to $Y=\omega^i\partial_{\theta_i}$, and let $g$ and
$\mathrm{II}$ be the induced first and second fundamental forms.
For the pressure-lifting Hamiltonian
\[
 H(\theta,I)=\omega\cdot I+\frac12 I^{\mathsf T}g^{-1}(\theta)I,
\]
a selected family of KAM graphs with
\[
 I=r d+o(r)
\]
gives exterior trace $X=Y+r g^{-1}d+o(r)$.  A sufficient condition for
positive pressure and a positive full surface energy term is
\[
 \omega\cdot d>0,
 \qquad \mathrm{II}(Y,g^{-1}d)<0\quad\hbox{everywhere on the interface}.
\]
Strictness and compactness give a uniform margin.  For sufficiently
small selected $r>0$,
\[
 p>0,\qquad
 C=\frac{\mathrm{II}(Y,Y)-\mathrm{II}(X,X)}{\mu_0}>0.
\]
Every term in the complete ideal-MHD quadratic energy is then
nonnegative, including the interface term.  This is a sufficient
linear energy-stability criterion; it does not assume the positivity
that was missing from the earlier calculation.

Its feasibility is explicitly one-dimensional.  Normalize
$\omega\cdot d=1$ and write
\[
 d=\frac{\omega}{|\omega|^2}+\sigma J\omega,
\]
where $J$ is rotation by $\pi/2$ in the two coordinate components.
Define
\[
 a_0(\theta)=\frac{\mathrm{II}(Y,g^{-1}\omega)}{|\omega|^2},
 \qquad b_0(\theta)=\mathrm{II}(Y,g^{-1}J\omega).
\]
The complete condition is simply
\[
 a_0(\theta)+\sigma b_0(\theta)<0\quad\hbox{for all }\theta.
\]
Equivalently, $\sigma$ must lie strictly above
$\sup_{b_0<0}(-a_0/b_0)$ and strictly below
$\inf_{b_0>0}(-a_0/b_0)$; every point with $b_0=0$ must have $a_0<0$.
Empty sets impose no bound.  Compactness supplies the uniform strict
margin once such a $\sigma$ exists.  Failure has a one-point or a
two-point witness.  This is a concrete geometric design test.

If noncollinearity is separately desired, one adds either sign of
\[
 \det(\omega,g^{-1}d)\ne0
\]
uniformly on the torus.  These are additional affine inequalities in
$\sigma$.  Positivity of $C$ itself does not imply noncollinearity;
noncollinearity is not needed merely to sign the displayed quadratic
energy when $C>0$ has already been proved.

\section{What remains before transferring the positive benchmark to a torus}

The cylindrical theorem removes a real difficulty: it gives an
interface gap uniform in the axial period, rather than a calculation
restricted to finitely many modes.  It does not, by itself, construct
the requested toroidal equilibrium.

There is an exact variational reduction that could support such a
transfer.  On a specified solid torus, the harmonic Neumann magnetic
field of fixed flux minimizes magnetic energy among all divergence-free
tangent fields with that flux.  If $B_h$ is this harmonic field, Hodge
orthogonality gives
\[
 \int |B|^2=\int|B_h|^2+\int|B-B_h|^2.
\]
At fixed total mass and fixed uniform specific entropy, Jensen's
inequality similarly minimizes the gas internal energy at uniform
density.  Thus the sum of the harmonic magnetic energies, at the
specified plasma and vacuum fluxes, and the uniform adiabatic gas
energy is an exact lower bound for the ideal potential energy on a
given plasma domain.  Its first shape derivative is the total-pressure
jump.  At the cylinder, the preceding mode calculation is its
interface Hessian, including the flux-dependent zero mode.

This suggests solving the Bernoulli shape equation by an elliptic
implicit-function theorem and preserving the positive Hessian.  A
bent tube of major radius $1/\varepsilon$ has, in cross-section and
arclength coordinates, the Euclidean metric
\[
 g_\varepsilon=dx^2+dy^2+(1+\varepsilon x)^2ds^2,
 \qquad L=\frac{2\pi}{\varepsilon}.
\]
For each fixed $L$, an ordinary local perturbation argument is not
enough: its permitted metric perturbation could depend on $L$.
The actual Euclidean closing condition ties $\varepsilon$ to $L$.
One must prove uniform estimates for the nonlinear harmonic-field
shape map and its inverse, including the magnetic-period constraints,
as $L$ increases.  Low axial Neumann modes make a naive estimate for
the scalar potentials singular.  Estimates for the fields and the
flux-constrained shape operator may cancel that singularity, but that
cancellation must be proved, not presumed from the interface spectrum.

A precise sufficient transfer certificate is a quadratic-form bound:
after identifying the boundaries and their flux constraints, the
change in the reduced interface Hessian must have operator norm less
than half the cylinder gap from $H^{1/2}$ to $H^{-1/2}$.  This would
preserve positivity of the whole energy because the bulk terms remain
nonnegative for a current-free core with constant pressure.  Such a
uniform bound and an exact toroidal Bernoulli solution have not been
established in this calculation.

There are further independent requirements.  A nonaxisymmetric
perturbation need not give a nonzero rotational transform, nested
magnetic surfaces, or exact quasisymmetry automatically.  The close
perfectly conducting wall and the edge current in the benchmark are
substantial assumptions.  Finite wall resistivity and kinetic wall
losses are absent.  Thus the positive result is a complete ideal-MHD
linear energy theorem for the stated benchmark and a concrete route
to a geometric extension, not a proof of hot stable fusion plasma.

\section{Constructive kinetic bounds and their precise
limits}\label{constructive-kinetic-bounds-and-their-precise-limits}

This section proves a second-order correction to the full-orbit
weak-quasisymmetry momentum and a finite-time, pointwise hot-core
theorem for linear Vlasov transport in a prescribed field with an
absorbing loss boundary. The theorem retains Maxwellian velocity tails
and permits a quantitative perturbation of the initial distribution. It
is not a result for self-consistent Vlasov--Maxwell--Landau. In
particular, it does not repair the equilibrium obstruction of the
original weak-QS seed. The new arbitrary harmonic vacuum-field
construction does not automatically satisfy weak QS and cannot be
substituted into the hypotheses without a separate proof.

\subsection{A corrected flux for full Lorentz
orbits}\label{a-corrected-flux-for-full-lorentz-orbits}

\paragraph{Hypotheses and
conventions}\label{hypotheses-and-conventions}

Let \(\Omega\subset\mathbb R^3\) be a bounded domain with \(C^2\)
boundary. Suppose \(B,u,\psi\) are \(C^3\) on a neighborhood of
\(\overline\Omega\), \(B\) never vanishes, and

\[
\nabla\cdot B=\nabla\cdot u=0,\qquad
B\times u=\nabla\psi,\qquad u\cdot\nabla |B|^2=0.
\tag{Q1}
\]

Write \(\beta=|B|,\ b=B/\beta,\ P=I-bb^T\), and

\[
S=\operatorname{sym}\nabla u,\qquad
K v=v\times b.
\]

Thus \(K^T=-K\) and \(K^2=-P\). All tensors and norms below are
Euclidean tensors and norms in physical space, not coordinate-component
norms in the embedded-torus chart.

Taking curl in \(B\times u=\nabla\psi\), and using both divergence
conditions, gives \([B,u]=0\). Consequently

\[
\operatorname{tr}S=0,\qquad
b^TSb=\frac{B\cdot((B\cdot\nabla)u)}{\beta^2}
     =\frac{u\cdot\nabla\beta^2}{2\beta^2}=0.
\tag{Q2}
\]

No assumptions about closed \(u\)-orbits, dense \(B\)-orbits, or
regularity of every interior flux level are required for the following
algebra.

\paragraph{Solving the gyrophase homological equation
explicitly}\label{solving-the-gyrophase-homological-equation-explicitly}

Put

\[
a=Sb,\qquad S_\perp=PSP,\qquad
C=b(Ka)^T+(Ka)b^T-\frac12S_\perp K,\qquad D=C/\beta.
\tag{Q3}
\]

Then \(a\perp b\), \(S_\perp\) has trace zero on \(b^\perp\), and

\[
C=C^T,\qquad [C,K]=CK-KC=S.
\tag{Q4}
\]

\textbf{Proof.} Decompose

\[
S=S_\perp+ba^T+ab^T.
\]

On the two-dimensional plane \(b^\perp\), a symmetric trace-zero matrix
anticommutes with the quarter-turn \(K\): \(S_\perp K+KS_\perp=0\). Thus
\(S_\perp K\) is symmetric. If \(C_1=b(Ka)^T+(Ka)b^T\), direct
multiplication using \(Kb=0\) and \(K^2a=-a\) gives
\([C_1,K]=ba^T+ab^T\). The same identities give
\([-\frac12S_\perp K,K]=S_\perp\). This proves both claims. In Frobenius
norm,

\[
\|C\|_F^2=2|a|^2+\frac14\|S_\perp\|_F^2
\le \|S\|_F^2.
\]

For any symmetric \(C\),

\[
\mathcal R_b(v^TCv)
:=(v\times b)\cdot\nabla_v(v^TCv)
=v^T(CK-KC)v.
\]

Therefore (Q4) is exactly the homological equation
\(\mathcal R_b(v^TCv)=v^TSv\). It is globally tensorial and requires no
choice of a gyroangle or frame. \(\square\)

\paragraph{Exact full-orbit
identity}\label{exact-full-orbit-identity}

For a species of mass \(m>0\) and charge \(q\ne0\), write the signed
quantity \(\varepsilon=m/q\). Define

\[
J(x,v)=\psi(x)-\varepsilon u(x)\cdot v
               +\varepsilon^2v^TD(x)v.
\tag{Q5}
\]

Consider prescribed-field Lorentz trajectories

\[
\dot x=v,\qquad
\dot v=\varepsilon^{-1}v\times B(x).
\tag{Q6}
\]

Then, exactly,

\[
\frac{d}{dt}J(x(t),v(t))
 =\varepsilon^2
   \sum_{i,j,k}v_i v_j v_k\,\partial_{x_k}D_{ij}(x(t)).
\tag{Q7}
\]

\textbf{Proof.} The first two terms of (Q5) give

\[
\frac d{dt}(\psi-\varepsilon u\cdot v)
 =v\cdot\nabla\psi-\varepsilon v^TSv-u\cdot(v\times B)
 =-\varepsilon v^TSv.
\]

For the last term,

\[
\frac d{dt}\big(\varepsilon^2v^TDv\big)
 =\varepsilon^2\sum_{ijk}v_iv_jv_k\partial_kD_{ij}
  +\varepsilon\beta\,\mathcal R_b(v^TDv).
\]

Since \(D=C/\beta\), the last summand equals \(\varepsilon v^TSv\) by
(Q4), and cancels the preceding defect. \(\square\)

This improves the exact first-order weak-QS momentum defect by one
order. It does not make \(J\) an exact invariant. The remainder is
explicitly computable from the physical-space derivatives of the seed.

\paragraph{A deterministic confinement
certificate}\label{a-deterministic-confinement-certificate}

Define

\[
U_0=\sup_{\overline\Omega}|u|,\qquad
D_0=\sup_{\overline\Omega}\|D\|_{\rm op},\qquad
D_1=\sup_{\substack{x\in\overline\Omega\\ |w|=1}}
\left|\sum_{ijk}w_iw_jw_k\partial_kD_{ij}(x)\right|.
\]

For \(V>0\), set

\[
A(V)=|\varepsilon|U_0V+\varepsilon^2D_0V^2,\qquad
R(V)=\varepsilon^2D_1V^3.
\]

For any trajectory segment remaining in \(\overline\Omega\), with
\(|v|=V_*\le V\),

\[
|\psi(x(t))-\psi(x(s))|
 \le 2A(V)+|t-s|R(V).
\tag{Q8}
\]

Indeed, \(|J-\psi|\le A(V)\), and integrating (Q7) gives
\(|J(t)-J(s)|\le |t-s|R(V)\). The speed is constant by (Q6). The same
argument works forwards and backwards in time.

Suppose in addition \(\psi=\psi_b\) on \(\partial\Omega\),
\(\psi<\psi_b\) in \(\Omega\), and define a core

\[
\Omega_\Delta=\{x\in\Omega:\psi(x)\le\psi_b-\Delta\}.
\]

Interior critical points of \(\psi\), including a magnetic axis, are
allowed. If

\[
\Delta>2A(V)+TR(V),
\tag{Q9}
\]

then every characteristic through \(x\in\Omega_\Delta,\ |v|\le V\), at
any time \(0\le t\le T\), remains inside \(\Omega\) when traced back to
time zero. For otherwise stop at the first backward boundary hit and
apply (Q8) on the segment before that hit. The change in \(\psi\) would
be at least \(\Delta\), contradicting (Q9). The strict inequality also
excludes a grazing boundary hit. The analogous forward no-hit assertion
holds for initial points in the core.

If \(R(V)>0\), this gives the explicit time

\[
T_{\rm cert}=\frac{\Delta-2A(V)}{R(V)}
\]

whenever its numerator is positive. If the numerator is nonpositive,
this certificate gives no information; its failure is not an escape
theorem. If \(R(V)=0\) and the numerator is positive, the algebraic
certificate holds for arbitrary time in the prescribed-field model.

\paragraph{Scaling}\label{scaling}

Under \(B\mapsto\Lambda B,\ \psi\mapsto\Lambda\psi,\ u\mapsto u\), one
has \(C\mapsto C,\ D\mapsto D/\Lambda\), while the geometric flux buffer
scales by \(\Lambda\). Consequently the cubic accumulated drift relative
to the flux buffer is of order \(T/\Lambda^2\). This is a precise
strong-field gain in this orbit certificate, not a statement about
self-consistent plasma stability at arbitrary field amplitude. No
numerical pass at the project's 10 T budget has been established here:
\(U_0,D_0,D_1\), and the physical flux buffer must first be bounded for
the actual embedded geometry.

\paragraph{Exact relativistic orbit
extension}\label{exact-relativistic-orbit-extension}

The orbit identity also has a direct relativistic version. Let \(w=p/m\), let
\(c_{\rm light}\) denote the speed of light, and set
\(\gamma(w)=(1+|w|^2/c_{\rm light}^2)^{1/2}\). For prescribed \(E=0\),
\[
\dot x=w/\gamma,\qquad
\dot w=(\varepsilon\gamma)^{-1}w\times B.
\] Both \(|w|\) and \(\gamma\) are constant. The vector field on
\((x,w)\) is exactly \(1/\gamma(w)\) times the Newtonian vector field
used in (Q6). Hence \[
J_{\rm rel}=\psi-\varepsilon u\cdot w+\varepsilon^2w^TDw,\qquad
\dot J_{\rm rel}=\frac{\varepsilon^2}{\gamma}
\sum_{ijk}w_iw_jw_k\partial_kD_{ij}.
\] On \(|w|\le W\), the endpoint bound is \(A(W)\), and the drift bound
is \[
R_{\rm rel}(W)=
\frac{\varepsilon^2D_1W^3}
{\sqrt{1+W^2/c_{\rm light}^2}}.
\] The function \(r^3/\sqrt{1+r^2/c_{\rm light}^2}\) is increasing for
\(r\ge0\), which justifies this uniform bound. The no-hit proof is
unchanged. The Maxwellian temperature theorem below remains explicitly
Newtonian; no identification of its scalar moment temperature with a
relativistic Jüttner rest-frame temperature is made.

\subsection{A pointwise hot-core theorem at an absorbing
boundary}\label{a-pointwise-hot-core-theorem-at-an-absorbing-boundary}

Consider the prescribed-field Vlasov equation

\[
\partial_t f+v\cdot\nabla_xf+
\varepsilon^{-1}(v\times B)\cdot\nabla_vf=0
\]

on \(\Omega\times\mathbb R^3\), with zero incoming data at the loss
boundary. Take

\[
f(0,x,v)=M(v)
=n_0(2\pi\sigma^2)^{-3/2}e^{-|v|^2/(2\sigma^2)},
\qquad \sigma^2=\Theta_0/m.
\tag{H1}
\]

Here \(\Theta_0\) is temperature in energy units. The initial trace and
zero incoming boundary data are not compatible at the time-zero corner;
the assertion is about the standard mild solution defined by backward
characteristics killed at their first boundary encounter, not a smooth
solution up to that corner.

Along a characteristic that reaches the initial surface, the value is
\(M(v(0))=M(v)\), since speed is conserved. Along one killed at the
boundary, the value is zero. Consequently

\[
0\le f(t,x,v)\le M(v)
\quad\hbox{everywhere},
\]

and (Q9) gives the stronger equality

\[
f(t,x,v)=M(v)
\quad\hbox{for }x\in\Omega_\Delta,\ |v|\le V,\ 0\le t\le T.
\tag{H2}
\]

No compact velocity support has been assumed. High-velocity particles
can be lost.

Define the local density, mean velocity and scalar kinetic temperature
by

\[
n=\int f\,dv,\quad U=n^{-1}\int vf\,dv,\quad
\Theta=\frac{m}{3n}\int |v-U|^2f\,dv.
\]

Let \(c=V/\sigma\), and define

\[
a(c)=\operatorname{erf}(c/\sqrt2)
       -\sqrt{2/\pi}\,c e^{-c^2/2},
\] \[
b(c)=\operatorname{erf}(c/\sqrt2)
       -\sqrt{2/\pi}\,(c+c^3/3)e^{-c^2/2},
\] \[
d(c)=\sqrt{2/\pi}(c^2+2)e^{-c^2/2}.
\]

Equivalently \(a=\mathbb P(\chi_3\le c)\) and
\(b=\mathbb P(\chi_5\le c)\). On the core and certified time interval,

\[
a(c)n_0\le n\le n_0,\qquad
\boxed{\quad
\frac{\Theta}{\Theta_0}
\ge b(c)-\frac{d(c)^2}{3a(c)^2}.
\quad}
\tag{H3}
\]

\textbf{Proof.} Integrating (H2) proves the lower density bound and
\(\int |v|^2f\ge3n_0\sigma^2b\); the global upper bound gives
\(n\le n_0\). The truncated Maxwellian has zero first moment. Thus the
entire mean-momentum contribution comes from the remaining tail, and

\[
\left|\int vf\,dv\right|
\le\int_{|v|>V}|v|M\,dv=n_0\sigma d.
\]

Using \(\Theta=m(\int |v|^2f/n-|\int vf|^2/n^2)/3\) and the density
bounds proves (H3). This explicitly allows the lost tail to be
asymmetric, and therefore does not assume zero mean flow after losses.
\(\square\)

For \(c=2.8\), the numerical constants are

\[
a>0.95056321,\quad
b-\frac{d^2}{3a^2}>0.82577125.
\]

Hence \(\Theta_0=12.5\,{\rm keV}\) implies \(\Theta>10.3221\,{\rm keV}\)
and \(n>0.95056n_0\) on the certified core and interval. For \(c=3\),
the temperature ratio is greater than \(0.88757292\). These strict
rational lower bounds are independently checked by the accompanying
alternating-series certificate. These are
implications of (Q9), not a claim that (Q9) has passed for the actual
device dimensions.

\paragraph{Stability under initial distribution
perturbations}\label{stability-under-initial-distribution-perturbations}

Suppose instead \[
(1-\delta)M\le f_0(x,v)\le(1+\delta)M,\qquad 0\le\delta<1.
\]

The same characteristic argument gives the same bounds for \(|v|\le V\)
in the core, and \(0\le f\le(1+\delta)M\) globally. If
\(d_0=2\sqrt{2/\pi}\), then

\[
n\ge(1-\delta)a n_0,\qquad
\frac{\Theta}{\Theta_0}\ge
\frac{1-\delta}{1+\delta}b
-\frac{(d+\delta d_0)^2}{3(1-\delta)^2a^2}.
\tag{H4}
\]

To prove the momentum bound used here, subtract the Maxwellian on the
protected ball. Its perturbation contributes at most
\(\delta n_0\sigma(d_0-d)\); the unprotected tail contributes at most
\((1+\delta)n_0\sigma d\), whose sum is \(n_0\sigma(d+\delta d_0)\). The
second-moment and density estimates are as before. At
\(c=2.8,\delta=0.01\), (H4) gives
\(\Theta/\Theta_0>0.80709315\), hence \(10.08866\,{\rm keV}\)
from \(12.5\,{\rm keV}\).

There is also a directional temperature bound, so this conclusion need
not hide a cold direction behind a large scalar trace. Let \[
\mathsf T=\frac mn\int(v-U)(v-U)^Tf\,dv.
\] The protected Maxwellian ball has second-moment matrix
\(n_0\sigma^2bI\). The same estimates therefore give \[
\frac{\mathsf T}{\Theta_0}\ \ge\
\left[\frac{1-\delta}{1+\delta}b
-\frac{(d+\delta d_0)^2}{(1-\delta)^2a^2}\right]I
\tag{H5}
\] as a quadratic-form inequality. For \(c=2.8,\delta=0\), its
coefficient is \(0.8078673232710732\), giving every directional
temperature above \(10.0983\,{\rm keV}\) from \(12.5\,{\rm keV}\). For
\(c=3,\delta=0.01\), its coefficient is \(0.8593547807238485\), giving
\(10.7419\,{\rm keV}\).

This is finite-time stability for an explicitly stated transport
problem. Electromagnetic feedback, Coulomb collisions, fusion, alpha
heating and fueling are absent from its hypotheses.

\subsection{Electric fields, collisions, and sources: the exact
additional
terms}\label{electric-fields-collisions-and-sources-the-exact-additional-terms}

Keep \(J\) constructed from the static reference weak-QS field \(B\).
Let the actual Lorentz field be \(B+\delta B(t,x)\), and include
\(E(t,x)\). Then

\[
\frac{dJ}{dt}
=\varepsilon^2\sum v_iv_jv_k\partial_kD_{ij}
 +(-u+2\varepsilon Dv)\cdot(E+v\times\delta B).
\tag{P1}
\]

This is an exact identity. In particular an electrostatic field constant
on flux surfaces need not annihilate the second-order term, even though
\(u\cdot E=0\).

The speed obeys \[
|v(s)|\le |v(t)|+|\varepsilon|^{-1}
\int_s^t\|E(r)\|_\infty\,dr
\] backwards (and the corresponding forward bound). For a known bound
\(V_*\) on the relevant speed range,

\[
|\dot J|\le
\varepsilon^2D_1V_*^3+
(U_0+2|\varepsilon|D_0V_*)
\big(\|E\|_\infty+V_*\|\delta B\|_\infty\big).
\tag{P2}
\]

Thus the orbit certificate survives with the integral of the right side
and \(A(V_*)\). A self-consistent theorem still has to prove these field
budgets from the initial data and actual sources; assuming them for all
time would assume the missing part of stability. With \(E\ne0\), the
Maxwellian bounds (H2)--(H4) also require accounting for the
initial-to-final speed change; they do not follow unchanged merely from
(P2).

For a prescribed kinetic collision generator in backward-observable form
\[
\mathcal L_v\phi=\ell(x,v)\cdot\nabla_v\phi
                 +a(x,v):D_v^2\phi,\qquad a\ge0,
\] one has \[
\mathcal L_vJ
=\ell\cdot(-\varepsilon u+2\varepsilon^2Dv)
 +2\varepsilon^2a:D.
\tag{P3}
\]

In an equivalent diffusion representation with covariance \(2a\), the
martingale part of \(J\) has quadratic-variation rate \[
2(-\varepsilon u+2\varepsilon^2Dv)^T
 a(-\varepsilon u+2\varepsilon^2Dv).
\tag{P4}
\]

Therefore collisions do not preserve the corrected invariant or the
deterministic protected-velocity ball. A stochastic first-exit estimate
needs the actual drift, diffusion, and high-speed tail bounds. The full
nonlinear Landau coefficients depend on the unknown distributions;
replacing them by fixed Maxwellian coefficients is a different model.

For completeness, there is a useful exact weak-form Landau estimate. For
one-species Coulomb Landau with kernel
\(a_C(w)=\kappa|w|^{-1}(I-\hat w\hat w^T)\), set \[
\mathcal D(f)=\frac12\iint ff'
(\nabla\log f-\nabla'\log f')^Ta_C(v-v')
(\nabla\log f-\nabla'\log f')\,dv\,dv'.
\] For a \(C^2\) observable \(\phi\) with
\(\|D_v^2\phi\|_{\rm op}\le H\), \[
\left|\int Q(f,f)\phi\,dv\right|
\le \sqrt{\kappa}\,H\,
n^{3/4}\left(\int|v|^2f\,dv\right)^{1/4}
\sqrt{\mathcal D(f)}.
\tag{P5}
\]

Indeed the weak form is minus one half of the pair integral of
\(ff'(\nabla\log f-\nabla'\log f')^Ta_C(\nabla\phi-\nabla'\phi')\).
Cauchy--Schwarz gives the square root of \(\mathcal D\) times the square
root of \(\frac12\iint ff'Y^Ta_CY\). The Lipschitz bound
\(|Y|\le H|v-v'|\) reduces the latter to
\((\kappa H^2/2)\iint ff'|v-v'|\), and
\(\iint ff'|v-v'|\le2n^{3/2}(\int|v|^2f)^{1/2}\).

For \(\phi=J\), the Hessian is \(2\varepsilon^2D\), so the collision
effect on the mean corrected flux is genuinely second order. First-order
collision contributions vanish because self-collisions conserve local
momentum. This is a mean-observable estimate, not an exit-probability or
temperature theorem. Spatial integration of (P5) still needs a density
bound (for example \(n\le N\)), or a substitute estimate. With that
bound, total mass \(M\), and second moment \(E_2\), \[
\left|\iint Q(f,f)\phi\,dv\,dx\right|
\le \sqrt{\kappa}\,H\sqrt N\,(ME_2)^{1/4}
       \sqrt{\int_\Omega\mathcal D(f)\,dx}.
\] It is essential not to declare \(N\) controlled for the nonlinear
open-boundary system without proving it. For multispecies collisions the
individual species momentum is not conserved; the mass-weighted sum is,
and the species-specific corrected fluxes do not automatically match
those weights.

For an added distribution source \(S_f\), the exact observable balance
acquires \(\iint J S_f\); outgoing absorbing particles also carry the
flux \(Jf(v\cdot n)\). Scalar heating power does not specify this source
moment or its entropy cost.

\subsection{Covariant derivatives and the remaining magnetic
shear}\label{covariant-derivatives-and-the-remaining-magnetic-shear}

For \(\mathcal R_a=(v\times a)\cdot\nabla_v\), one has \[
[\mathcal R_a,\mathcal R_c]=\mathcal R_{a\times c}.
\] Define \[
a_i=(\partial_i b)\times b,\qquad
\mathfrak D_i=\partial_{x_i}+\mathcal R_{a_i},\qquad
\Omega_c=q\beta/m.
\] The vector identity \(a_i\times b=-\partial_i b\) gives \[
[\mathfrak D_i,\Omega_c\mathcal R_b]
 =(\partial_i\Omega_c)\mathcal R_b.
\tag{C1}
\] Thus the derivative of the \textbf{direction} of the magnetic field
cancels exactly. This removes its factor \(\beta\). It leaves the
derivative of the cyclotron frequency.

The remaining exact transport commutators are \[
[\mathfrak D_i,v\cdot\nabla_x]
 =(v\times a_i)\cdot\nabla_x
   -\mathcal R_{(v\cdot\nabla_x)a_i},
\tag{C2}
\] \[
[\mathfrak D_i,(q/m)E\cdot\nabla_v]
 =(q/m)(\partial_iE+a_i\times E)\cdot\nabla_v.
\tag{C3}
\] In (C2), the last rotation parameter is itself linear in \(v\), so
its coefficient before \(\nabla_v\) is quadratic in velocity. The
reduction of a large field-direction coefficient is useful, but it does
not eliminate geometric derivatives or velocity-weight requirements.

For this particular minimal-rotation connection, \[
[\mathfrak D_i,\mathfrak D_j]
 =-[\,b\cdot(\partial_i b\times\partial_j b)\,]\mathcal R_b.
\tag{C4}
\] To check (C4), differentiate \(a_i=\partial_i b\times b\); the
derivative part is \(-2\partial_i b\times\partial_j b\), and
\(a_i\times a_j=\partial_i b\times\partial_j b\). The cross product is
parallel to \(b\). One can modify the connection by adding
\(\alpha_i b\). A full chosen orthonormal frame can give a flat
connection locally; thus (C4) should not be advertised as an
unconditional topological obstruction. The corresponding streaming
connection terms remain, and (C1)'s magnitude derivative is unaffected
by such a frame choice.

Rotational covariance of the Coulomb kernel gives \[
\mathfrak D_iQ(f,g)=Q(\mathfrak D_i f,g)+Q(f,\mathfrak D_i g).
\tag{C5}
\] For a spatially uniform isotropic Maxwellian reference, the
linearized Landau operator therefore commutes with \(\mathfrak D_i\). It
does not follow that a linearization around a spatially varying
anisotropic reference commutes.

The remaining magnitude shear cannot be removed by a smooth periodic
gyroangle gauge. At fixed speed and pitch, write the rotation as
\(\Omega_c(x)\partial_\vartheta\). A derivative
\(\partial_i+\alpha(x,\vartheta)\partial_\vartheta\) would commute with
it only if \[
\partial_\vartheta\alpha=\partial_i\log|\Omega_c|.
\] Integrating around the gyroangle circle forces
\(\partial_i\Omega_c=0\). The pure rotation subproblem already exhibits
the shear: \[
g(t,x,\vartheta)=g_0(x,\vartheta-\Omega_c(x)t)
\] has a spatial derivative containing
\(-t(\partial_i\Omega_c)\partial_\vartheta g_0\). This is an obstruction
to that proposed cancellation, not a theorem that the full collisional
transport system is unstable. A time-dependent transported derivative
can cancel the shear, but its coefficients grow in time and a uniformly
equivalent norm still needs proof.

\paragraph{A global rotating frame on the solid
torus}\label{a-global-rotating-frame-on-the-solid-torus}

On a solid torus the oriented two-plane bundle \(b^\perp\) is trivial:
oriented two-plane bundles are classified by their Euler class in
\(H^2(\Omega;\mathbb Z)=0\). Thus a smooth global \(R(x)\in SO(3)\) with
\(R(x)e_3=b(x)\) exists. Put \[
A_i=R^T\partial_iR,\qquad g(x,w)=f(x,Rw).
\] Each \(A_i\) is skew. The prescribed transport becomes exactly \[
\partial_tg+(Rw)\cdot\nabla_xg+
\left[\Omega_c(w\times e_3)
      -\sum_i(Rw)_iA_iw\right]\cdot\nabla_wg.
\tag{C6}
\] This follows by the chain rule:
\(\partial_i f=\partial_i g-(A_iw)\cdot\nabla_wg\) at fixed physical
velocity. The physical covariant derivatives corresponding to fixed
\(w\) form a flat pure-gauge connection. Their rotation parameters equal
\(\partial_i b\times b+\alpha_i b\) for suitable \(\alpha_i\), so they
still commute with \(\mathcal R_b\).

The streaming and connection terms in (C6) satisfy the exact divergence
cancellation \[
\nabla_x\cdot(Rw)
+\nabla_w\cdot\left(-\sum_i(Rw)_iA_iw\right)=0.
\tag{C7}
\] Indeed the first term is \(\sum_i(R^Te_i)^TA_iw\), and the second is
its negative because \(\operatorname{tr}A_i=0\). Also \(w\cdot A_iw=0\).
Every radial velocity weight is therefore invariant under the connection
transport, and this transport is skew in weighted \(L^2(dx\,dw)\), up to
the actual spatial boundary flux. The Jacobian of
\((x,w)\mapsto(x,R(x)w)\) is one. Landau's rotational covariance
preserves its collision form.

Under \(B\mapsto\Lambda B\), \(R,A_i\) remain unchanged. Only
\(\Omega_c\) is multiplied by \(\Lambda\). Differentiation still
produces \((\partial_i\Omega_c)\mathcal R_{e_3}g\), derivatives of \(R\)
multiplying \(|w|\nabla_xg\), and derivatives of \(A_i\) multiplying
\(|w|^2\nabla_wg\). This global-frame formulation removes the
field-direction large parameter at zero order and displays where
weighted high-order estimates must still close.

\paragraph{A non-tautological weighted estimate for a simpler
collision
operator}\label{a-non-tautological-weighted-estimate-for-a-simpler-collision-operator}

There is a useful constructive check that the remaining frequency shear
need not imply instability. Freeze spatial transport and use the
normalized Ornstein--Uhlenbeck generator \(L=\Delta_w-w\cdot\nabla_w\)
on Gaussian \(L^2(M\,dw)\): \[
\partial_tg+\Omega_c(x)\mathcal R_{e_3}g=\nu Lg.
\] For one spatial derivative put \(h=\partial_i g\),
\(a(x)=1+|\partial_i\Omega_c|^2/\nu^2\), and \[
\mathcal E_i=\int\big(|h|^2+a|\mathcal R_{e_3}g|^2\big)M\,dw\,dx.
\] Commutation of \(L\) with rotations, Gaussian integration by parts,
and Gaussian Poincare give \[
\frac d{dt}\mathcal E_i
\le-\nu\int|\nabla_wh|^2M
    -\nu\int(a+1)|\nabla_w\mathcal R_{e_3}g|^2M.
\tag{C8}
\] To verify this, the only cross term is
\(-2(\partial_i\Omega_c)\langle h,\mathcal R_{e_3}g\rangle_M\). The
rotation has zero Gaussian mean. Subtract the Gaussian mean of \(h\),
apply Poincare and Young, and then apply Poincare to
\(\mathcal R_{e_3}g\). The cross term costs at most
\(\nu\|\nabla_wh\|^2+\nu(a-1)\|\nabla_w\mathcal R_{e_3}g\|^2\), which
yields (C8).

This is an actual modified energy estimate, but only for the stated
subproblem. Restoring spatial streaming differentiates \(a(x)\) and
introduces the connection terms in (C6). Full Landau has momentum and
energy null modes and weak high-velocity coercivity, so the two Gaussian
Poincare substitutions do not give (C8) for that operator. Those are
concrete missing estimates, not an inference that no such estimate can
exist.

For a vacuum field, \(\nabla\cdot B=\nabla\times B=0\) implies
\(\Delta B=0\), hence \[
\Delta |B|^2=2|\nabla B|^2.
\] Consequently a vacuum field with constant magnitude is constant as a
vector field. A nontrivial vacuum stellarator cannot generally evade
(C1) by choosing a nonconstant field of constant magnitude.

\subsection{What remains unproved}\label{what-remains-unproved}

The finite-time transport theorem is a genuine consequence of explicit
geometric bounds, provided its buffer inequality is satisfied. We have
not certified that inequality for the project's actual embedded geometry
and field budget. The common-temperature Maxwellian used here has no
self-consistent field response, and the prescribed original weak-QS
field has a separately established obstruction to the proposed
equilibrium classes. Thus this construction is an orbit-level and
kinetic-transport advance, not an externally coiled hot fusion
equilibrium. Extending it requires actual collision and source
estimates, electromagnetic feedback bounds, wall physics and a globally
coercive continuation estimate. The covariant calculation removes the
large field-direction commutator but leaves field-strength shear,
spatial transport terms and velocity-weight requirements.

\subsection{Primary sources and method
provenance}\label{primary-sources-and-method-provenance}

The explicit corrected flux and absorbing-core inequalities above are calculations made here. The following primary references delimit adjacent established results; the geometric inspiration is discussed earlier in this note.

Karaki's arbitrary-smooth-field hypocoercivity result concerns a linear
Fokker--Planck operator on the flat torus:

\url{https://doi.org/10.3934/krm.2020011}

Fan--Lei's full Vlasov--Maxwell--Landau result allows a strong uniform
background magnetic field:

\url{https://doi.org/10.1007/s10955-022-03030-1}

Dong--Guo--Yastrzhembskiy's global bounded-domain relativistic
Vlasov--Maxwell--Landau theorem is near Jüttner, with electromagnetic
smallness and zero magnetic flux through toroidal cuts:

\url{https://arxiv.org/abs/2401.00554}

Ouyang--Wu's October 2026 preprint gives a global near-Maxwellian
theorem for the Landau equation with inflow or diffuse reflection, but
does not assert a full strong-variable-field Vlasov--Maxwell--Landau
theorem:

\url{https://arxiv.org/abs/2610.40517}

Carrapatoso--Mischler give a near-equilibrium Landau theorem with
Maxwell accommodation and quantitative convergence:

\url{https://arxiv.org/abs/2407.09031}


\section{A concrete parameter illustration and the next mathematical target}

Within the cylindrical ideal-MHD theorem only, take equal electron and ion number densities $n_e=n_i=10^{20}\,\mathrm{m}^{-3}$ and temperatures in energy units $\Theta_e=\Theta_i=12.5\,\mathrm{keV}$. Their ideal-gas pressure is
\[
p=n_e\Theta_e+n_i\Theta_i\simeq4.00544\times10^5\,\mathrm{Pa}.
\]
Set $B_i=5\,\mathrm T$ and define the dimensionless parameter by the exact relation $s=8\mu_0p/B_i^2$. The resulting values are approximately
\[
s\simeq0.16107,\quad \beta_{\rm plasma}=s/4\simeq0.04027,
\quad B_e\simeq4.6883\,\mathrm T,\quad B_\theta\simeq2.0067\,\mathrm T.
\]
The theorem allows any wall satisfying
\[
1<b/a\le1+s/16\simeq1.01007.
\]
The proof uses the exact symbolic relations and inequalities, not these rounded decimal displays. This is a consistent hot ideal-MHD benchmark with a close ideal wall and a substantial interface current. It is not a parameter certificate for an embedded stellarator, a cold absorbing reactor wall, kinetic stability, or net fusion power.

There are now two concrete geometric ways to continue. One is to search the exact harmonic family for a toroidal interface satisfying the explicit curvature feasibility inequalities, or to bound the full boundary operator directly when those inequalities fail. The other is to solve the toroidal Bernoulli shape equation by continuing the positive cylindrical branch, with estimates uniform in the axial period. The latter branch is particularly informative: its transverse field is first order in a small amplitude, its axial-field decrease is second order, and its pressure is positive at second order. The resonant counterexample shows why a generic positive first-order pressure shift is not a stability argument.

Neither geometric route completes the kinetic problem. The same final configuration must also admit a self-consistent distribution, collision and wall-loss estimates, bounded electromagnetic perturbations, a specified heating/fueling/ash model, and an actual coil realization with certified errors. The corrected invariant and covariant derivatives provide terms to estimate in such a proof; they do not permit these requirements to be assumed.

\section{Reproducibility contents}

The accompanying source bundle contains this note's LaTeX sources, the original rational candidate data, the converted 121-monomial Laurent potential, the exact symbolic verification script, the rational Maxwellian-bound certificate, and the pinned source-dependency audit. The seed potential is specified completely by \texttt{dommaschk\_complex\_laurent.json}: each record contributes $(a+ib)u^p v^q z^j$ with exact rational strings. Its physical field is the Euclidean restriction of
\[
\frac{1}{2i}\left(\frac{du}{u}-\frac{dv}{v}\right)+dF_0.
\]
The field-family formula and basis theorem in this note specify every perturbation without numerical fitting. Running \texttt{complex\_vacuum\_exact\_check.py} with SymPy reproduces the vacuum and symmetry identities. Running \texttt{verify\_hot\_core\_constants.py} requires only the Python standard library and verifies rational lower bounds using alternating-series enclosures and integer square roots. Neither script claims to certify KAM radii, the actual weak-QS loss time, a toroidal free-boundary spectrum, or a nonlinear plasma evolution.

\begin{thebibliography}{99}
\bibitem{wave} OpenAI, \emph{Local smoothing for the Euclidean wave equation}, 25 October 2026, Sections 2--4 and Proposition 4.2. Primary manuscript: \url{https://github.com/openai/math/tree/main/preprints/Local-smoothing-for-the-Euclidean-wave-equation-October-25-2026}.
\bibitem{dlL} R. de la Llave, A. González, A. Jorba and J. Villanueva, \emph{KAM theory without action-angle variables}, Nonlinearity 18 (2005), 855--895, Theorem 1 and Definition 2. \url{https://doi.org/10.1088/0951-7715/18/2/020}.
\bibitem{hitchin} N. Hitchin, \emph{Monopoles and geodesics}, Communications in Mathematical Physics 83 (1982), 579--602. \url{https://doi.org/10.1007/BF01208717}.
\bibitem{walsh} J. L. Walsh, \emph{The approximation of harmonic functions by harmonic polynomials and by harmonic rational functions}, Bulletin of the AMS 35 (1929), 499--544. \url{https://doi.org/10.1090/S0002-9904-1929-04753-0}.
\bibitem{elps} A. Enciso, A. Luque and D. Peralta-Salas, \emph{MHD equilibria with nonconstant pressure in nondegenerate toroidal domains}, Journal of the European Mathematical Society. \url{https://doi.org/10.4171/JEMS/1410}.
\bibitem{plasmalean} PlasmaLean project, source commit \texttt{6312f1cc9c30b32798ad7b7540766152ef03e76f}, especially \texttt{DommaschkKamAFinal}, \texttt{DommaschkKamCFinal}, \texttt{DommaschkKamNested}, and the admissibility/frame modules. \url{https://github.com/siddhartha-gadgil/PlasmaLean}.
\end{thebibliography}
\end{document}
